% GENERATED FILE. Edit canonical lessons, then run npm run generate:books.
% canonical-source-hash: fnv1a64:26d1153a9f29ff22
% canonical-lessons: TE-C252-ugadi, TE-C252-dipavali, TE-C252-dasara, TE-C252-vinayaka-caviti, TE-C252-batukamma

\chapter{Ugadi, Diwali, Dasara, Vinayaka Chavithi, Bathukamma}
\label{ch:te-ugadi-diwali-dasara-vinayaka-chavithi-bathukamma}
\hlchaptermodality{252}

\begin{chapteropening}
\textbf{By the end of this chapter:} \emph{I can say Ugadi, Diwali, Dasara, Vinayaka Chavithi and Bathukamma.}

Complete the last lesson of chapter 252: I can say Ugadi, Diwali, Dasara, Vinayaka Chavithi and Bathukamma.
\end{chapteropening}

\section[\texorpdfstring{\te{ఉగాది} (ugādi)}{ఉగాది (ugādi)}]{\te{ఉగాది} (ugādi) — Ugadi, the Telugu New Year}

\label{lesson:TE-C252-ugadi}



\begin{quote}
Before the new one: say the Telugu for an injection, then the Telugu for a vaccination.
\end{quote}

\subsection*{You'll want to know: \te{ఉగాది}}
\textbf{\te{ఉగాది}} — \emph{ugādi} — \textquotedblleft{}Ugadi, the Telugu New Year\textquotedblright{}.

\begin{grammarlens}[title={What you've built}]
One more word for this chapter.
\end{grammarlens}

\subsection*{Your turn}
\begin{itemize}
\raggedright
  \item \emph{Say it:} \emph{ugādi}
  \item \emph{Say it:} \emph{ugādi}, once more
  \item \emph{Say it:} say \emph{ugādi}
  \item \emph{Recall:} say the Telugu for an injection, then the Telugu for a vaccination, then say \emph{ugādi} again
\end{itemize}

\begin{tcolorbox}[breakable,colback=teal!4,colframe=teal!35!black,title={Before you move on}]
What does \te{ఉగాది} mean? (\textquotedblleft{}Ugadi, the Telugu New Year\textquotedblright{}.) Say it once more.
\end{tcolorbox}
\section[\texorpdfstring{\te{దీపావళి} (dīpāvaḷi)}{దీపావళి (dīpāvaḷi)}]{\te{దీపావళి} (dīpāvaḷi) — Diwali, the festival of lamps}

\label{lesson:TE-C252-dipavali}



\begin{quote}
Before the new one: say the Telugu for a vaccination, then the Telugu for Ugadi.
\end{quote}

\subsection*{You'll want to know: \te{దీపావళి}}
\textbf{\te{దీపావళి}} — \emph{dīpāvaḷi} — \textquotedblleft{}Diwali, the festival of lamps\textquotedblright{}.

\begin{grammarlens}[title={What you've built}]
One more word for this chapter.
\end{grammarlens}

\subsection*{Your turn}
\begin{itemize}
\raggedright
  \item \emph{Say it:} \emph{dīpāvaḷi}
  \item \emph{Say it:} \emph{dīpāvaḷi}, once more
  \item \emph{Say it:} say \emph{dīpāvaḷi}
  \item \emph{Recall:} say the Telugu for a vaccination, then the Telugu for Ugadi, then say \emph{dīpāvaḷi} again
\end{itemize}

\begin{tcolorbox}[breakable,colback=teal!4,colframe=teal!35!black,title={Before you move on}]
What does \te{దీపావళి} mean? (\textquotedblleft{}Diwali, the festival of lamps\textquotedblright{}.) Say it once more.
\end{tcolorbox}
\section[\texorpdfstring{\te{దసరా} (dasarā)}{దసరా (dasarā)}]{\te{దసరా} (dasarā) — Dasara}

\label{lesson:TE-C252-dasara}



\begin{quote}
Before the new one: say the Telugu for Ugadi, then the Telugu for Diwali.
\end{quote}

\subsection*{You'll want to know: \te{దసరా}}
\textbf{\te{దసరా}} — \emph{dasarā} — \textquotedblleft{}Dasara\textquotedblright{}.

\begin{grammarlens}[title={What you've built}]
One more word for this chapter.
\end{grammarlens}

\subsection*{Your turn}
\begin{itemize}
\raggedright
  \item \emph{Say it:} \emph{dasarā}
  \item \emph{Say it:} \emph{dasarā}, once more
  \item \emph{Say it:} say \emph{dasarā}
  \item \emph{Recall:} say the Telugu for Ugadi, then the Telugu for Diwali, then say \emph{dasarā} again
\end{itemize}

\begin{tcolorbox}[breakable,colback=teal!4,colframe=teal!35!black,title={Before you move on}]
What does \te{దసరా} mean? (\textquotedblleft{}Dasara\textquotedblright{}.) Say it once more.
\end{tcolorbox}
\section[\texorpdfstring{\te{వినాయక} \te{చవితి} (vināyaka caviti)}{వినాయక చవితి (vināyaka caviti)}]{\te{వినాయక} \te{చవితి} (vināyaka caviti) — Vinayaka Chavithi, Ganesha's festival}

\label{lesson:TE-C252-vinayaka-caviti}



\begin{quote}
Before the new one: say the Telugu for Diwali, then the Telugu for Dasara.
\end{quote}

\subsection*{You'll want to know: \te{వినాయక} \te{చవితి}}
\textbf{\te{వినాయక} \te{చవితి}} — \emph{vināyaka caviti} — \textquotedblleft{}Vinayaka Chavithi, Ganesha's festival\textquotedblright{}.

\begin{grammarlens}[title={What you've built}]
One more word for this chapter.
\end{grammarlens}

\subsection*{Your turn}
\begin{itemize}
\raggedright
  \item \emph{Say it:} \emph{vināyaka caviti}
  \item \emph{Say it:} \emph{vināyaka caviti}, once more
  \item \emph{Say it:} say \emph{vināyaka caviti}
  \item \emph{Recall:} say the Telugu for Diwali, then the Telugu for Dasara, then say \emph{vināyaka caviti} again
\end{itemize}

\begin{tcolorbox}[breakable,colback=teal!4,colframe=teal!35!black,title={Before you move on}]
What does \te{వినాయక} \te{చవితి} mean? (\textquotedblleft{}Vinayaka Chavithi, Ganesha's festival\textquotedblright{}.) Say it once more.
\end{tcolorbox}
\section[\texorpdfstring{\te{బతుకమ్మ} (batukamma)}{బతుకమ్మ (batukamma)}]{\te{బతుకమ్మ} (batukamma) — Bathukamma, the Telangana flower festival}

\label{lesson:TE-C252-batukamma}



\begin{quote}
Before the new one: say the Telugu for Dasara, then the Telugu for Vinayaka Chavithi.
\end{quote}

\subsection*{You'll want to know: \te{బతుకమ్మ}}
\textbf{\te{బతుకమ్మ}} — \emph{batukamma} — \textquotedblleft{}Bathukamma, the Telangana flower festival\textquotedblright{}.

\begin{grammarlens}[title={What you've built}]
One more word for this chapter.
\end{grammarlens}

\subsection*{Your turn}
\begin{itemize}
\raggedright
  \item \emph{Say it:} \emph{batukamma}
  \item \emph{Say it:} \emph{batukamma}, once more
  \item \emph{Say it:} say \emph{batukamma}
  \item \emph{Recall:} say the Telugu for Dasara, then the Telugu for Vinayaka Chavithi, then say \emph{batukamma} again
\end{itemize}

\begin{tcolorbox}[breakable,colback=teal!4,colframe=teal!35!black,title={Before you move on}]
What does \te{బతుకమ్మ} mean? (\textquotedblleft{}Bathukamma, the Telangana flower festival\textquotedblright{}.) Say it once more.
\end{tcolorbox}
